\scanpage{8}

\deleted{On peut faire la construction de $K(\bar C,\bar\delta)$ en trois étapes $K(C,\delta)=\rho_2\rho_1\rho_0(C,\delta)$, avec \unsure{sa|par} construction transfinie comme précédemment, en ajoutant les conditions \unsure{$\beta$ e)|$\rho_2$)} aux conditions a) b) c) d)\\ e) \unsure{$e_0$)|$\rho_0$)} Si $R\xrightarrow{\varphi}R'$ est un des $f'_\lambda$, et $u,v:R'\rightrightarrows C_i$.}

Utilisant le th. précédent, on peut supposer que $C\in\mathrm{Ob}\,\Delta$, et $\delta'$ \inserted{$(\delta'')$} contient les \deleted{diagrammes}\inserted{cônes pr.}\inserted{ (ind.)} exacts de type $d'$ ($d''$), de sorte que
\[
\underline{\mathrm{Hom}}_\delta(C,E)=\underline{\mathrm{Hom}}_{\Delta,\delta}(C,E)\qquad\forall E\in\mathrm{Ob}\,\Delta.
\]

On va faire une construction transfinie, $(C_i,\delta_i)$, $C_i$ \inserted{petite} $\in\mathrm{Ob}\,\Delta$ foncteurs de transition dans $\Delta$,
\[
\left\{\begin{aligned}
&(C_0,\delta_0)=(C,\delta)\\
&(C_{i+1},\delta_{i+1})=K(C_i,\delta_i),\quad\tau_i:C_i\unsure{-|\to}C_{i+1},\quad\delta_{i+1}=\tau_i(\delta_i)\\
&C_i=\varinjlim_{j<i\,(\Delta)}C_j\text{ si }i\text{ ordinal limite},\quad\delta_i=\bigcup_{\underline j<i}\text{Images des }\delta_j\text{ dans }C_i.\ \deleted{\unsure{\text{categ.}}}
\end{aligned}\right.
\]

On construit $(C,\delta)\longmapsto K(C,\delta)$ de façon à satisfaire les conditions suivantes

a) \quad $a_0$) Pour tt foncteur $R\xrightarrow{\varphi}R'$ parmi les $f^{(i)}_\lambda$ \inserted{$(i=0,1,2)$}, et pour tout couple de \inserted{$\Delta$-}foncteurs $u',v':R'\rightrightarrows C_i$, et une double flèche $\lambda,\mu:u'-v'$, si $\lambda*\varphi=\mu*\varphi$, alors $\tau_i(\lambda)=\tau_i(\mu)$

$a_1$) Pour tte \deleted{\unsure{$\Delta$- et tte flèche, soit}} $\varphi:R-R'$ parmi les $f^1_\lambda$ et $f^2_\lambda$, et tt couple de $\Delta$-foncteurs $u',v':R'-C_i$, et $\lambda':u'\circ\varphi\to v'\circ\varphi$, \deleted{\unsure{$\exists$ un}} considérant $\tau_i\circ u'\xrightarrow{\tau_i*\lambda'}\tau_i\circ v'$,

\diagram{(\begin{tikzcd}[ampersand replacement=\&,column sep=large] R^{\unsure{\prime|{}}}\arrow[r,shift left=1ex,"\tau_i u'\varphi"{name=U}]\arrow[r,shift right=1ex,swap,"\tau_i v'\varphi"{name=V}]\& C_{i+1}\arrow[from=U,to=V,shorten <=1pt,shorten >=1pt]\end{tikzcd})}

\deleted{\unsure{$\exists$}} $\exists\,\unsure{\bar\lambda|\overset{\text{un}}{\lambda}|\lambda}:\tau_i\circ u\deleted{\circ\varphi}\longrightarrow\tau_i\circ v\deleted{\circ\varphi}$ tel que $\tau_i*\lambda'\deleted{*\varphi}=\deleted{\unsure{\varphi*}}\bar\lambda*\varphi$

$a_2$) Pour tt $\varphi:R\longrightarrow R'$ parmi les $f^2_\lambda$, et tt $\Delta$-foncteur $u:R\longrightarrow C_i$, \deleted{$\exists$} d'où $\tau_i\unsure{:|\circ}u:R-C_{i+1}$, $\exists$ un $\Delta$-foncteur $u':\deleted{\unsure{\Delta}}R'\longrightarrow C_{i+1}$ et un isom
\[
\tau_i\circ u\simeq u'\circ\varphi
\]

% archive sheet number: 7
